\documentclass[11pt]{article}
\usepackage[margin=1in]{geometry}
\usepackage{amsmath,amssymb}
\usepackage{booktabs}
\usepackage{hyperref}
\usepackage{xcolor}

\title{First-Pass Replication Report: Friedrich, Durante, Scholz (2012)\\
\large Static GLOBLE / Giant-Loop DNA Damage Model (GLDM) --- DSB clustering in Mbp chromatin loops}
\author{LUCID Replication Slot 63 (Backfilled 2026-07-06)}
\date{First pass 2026-06; audit 2026-06-20; backfill 2026-07-06}

\begin{document}
\maketitle

\section*{Verdict}

\textbf{PARTIAL.} The static GLOBLE analytic model (paper Eqs.~1--7) was
reimplemented from scratch in \verb|code/globle_static.py| using the paper's
abstract plus the well-documented kinetic-GLOBLE extension (Herr et al.\ 2014,
PLoS ONE) as a formula cross-check. The reimplementation reproduces the
paper's central dose-response shape claim (LQ at low $D$, quasi-linear at
intermediate clinical dose) and the analytic map from microscopic parameters
$(\varepsilon_i, \varepsilon_c, \alpha_{DSB}, N_L)$ to LQ coefficients
$(\alpha_{LQ}, \beta_{LQ})$. The cross-cell-line $\beta$-vs-$\alpha$
anti-correlation --- one of the paper's headline claims --- was
\emph{not} confirmed: on the 17-cell-line subset available from Herr 2014
Table~2 we measured Pearson $r(\alpha,\beta)=+0.655$ (opposite sign).
The paper's own claim was anchored on a 150+ cell-line PIDE meta-analysis
that we did not re-digitise. External 3-judge audit (2026-06-20)
independently returned PARTIAL (median Coverage 4/10, Agreement 6/10).

\section{What the paper's headline actually is}

Friedrich, Durante, and Scholz (RR2964, 2012) is the founding static
formulation of the GLOBLE / Giant-Loop DNA Damage Model (GLDM) family. Its
headline contributions are:

\begin{enumerate}
\item Cell survival after photon irradiation is derived from a
\emph{two-lesion-class} statistic: DSBs that fall alone within one of the
$N_L \approx 3000$ megabase-pair chromatin loops are ``isolated'' (mild
lethality $\varepsilon_i$); loops containing $\geq 2$ DSBs are ``clustered''
(large lethality $\varepsilon_c$). Under Poisson placement of DSBs into loops,
$-\ln S(D)$ is a closed-form sum of two terms.
\item The model naturally predicts the empirical LQ-to-linear crossover
observed in high-dose photon survival curves (paper Fig.~1--3).
\item The paper's most-cited empirical claim: across cell lines, the
LQ coefficients $\alpha$ and $\beta$ are \emph{anti-correlated}, and this
anti-correlation is a direct consequence of $\beta_{LQ}$ having a
$-2\varepsilon_i$ contribution when $\varepsilon_c$ is roughly bounded
(paper Fig.~5 / Eq.~12--13). The claim was anchored on a 150+ cell-line
compilation (the future PIDE database).
\item Downstream: the same static machinery is later extended to ions and
mixed radiation quality (Friedrich et al.\ 2013 IJRB; Herr et al.\ 2014
kinetic GLOBLE), giving the whole GLDM/GLOBLE line its predictive power.
\end{enumerate}

\section{What we actually did (headline exercised: PARTIAL)}

\begin{itemize}
\item \textbf{Analytic model:} Yes --- fully reimplemented. Eqs.~1--7 for
$\lambda(D)$, $p_0/p_i/p_c$, $n_i/n_c$, $-\ln S$, plus the LQ correspondence
Eqs.~12--13. Verified numerically at $D \in \{2,5,10,20\}$~Gy.
\item \textbf{Dose-response shape (Claim~1):} Yes --- reproduced. Fig.~1
(RT112 dose-response) and Fig.~3 (isolated/clustered decomposition) both
show the LQ-tangent-to-linearising transition the paper highlights, driven
mechanically by the $p_c$ term.
\item \textbf{LQ mapping (Claim~3):} Yes --- reproduced analytically and
numerically. RT112 gives $\alpha \approx 0.159$/Gy, $\beta \approx 0.0277$/Gy$^2$,
$\alpha/\beta \approx 5.7$~Gy, in the literature ballpark.
\item \textbf{$\beta$-vs-$\alpha$ anti-correlation (Claim~2):}
\textbf{NOT confirmed on our subset.} On the 17 lines available
(C3H~10T1/2, CHO~10B2, CHO~K1, NFF28, HX118, HX32, HX58, MT, LL, B16, HX34,
IN859, IN1265, SB, RT112, HX138, HX142) we measured Pearson $r=+0.655$,
Spearman $\rho=+0.512$ --- opposite sign to the paper. Mechanistically:
in this small sample $\varepsilon_i$ and $\varepsilon_c$ are themselves
strongly positively correlated, which masks the paper's predicted
$\beta$-vs-$\alpha$ anti-correlation.
\end{itemize}

\section{What we did NOT do (honest scope)}

\begin{enumerate}
\item \textbf{No track-structure Monte Carlo.} The paper's DSB yield
$\alpha_{DSB} = 30$~DSB/Gy/cell is taken as a paper-fixed constant; we did
not re-derive it from PARTRAC or any nanoscale-MC track-structure code.
For photons this is a defensible short-cut (the constant is well-established
and roughly LET-independent in the photon regime), but it means our
replication does \emph{not} exercise the ``track structure $\rightarrow$
DSB spectrum'' half of the GLDM story.
\item \textbf{No DSB-clustering algorithm reimplementation on real chromatin
geometries.} The paper's analytic Poisson-placement into loops
(Eqs.~4--7) is what we reimplemented. We did \emph{not} build a
DSB-cluster-detection algorithm on any explicit chromatin model
(fractal-globule, Hi-C-informed loop domains, polymer simulation, etc.)
and rerun cluster spectra from geometry. The paper itself relies on the
Poisson approximation over a mean loop-length distribution; a modern
Hi-C-anchored geometric rerun would be a real extension, not a
replication.
\item \textbf{No PDF text access.} The paper is closed-access
(Unpaywall \verb|is_oa=false|). We worked from PubMed abstract + Herr~2014
sibling paper. Risk of minor symbol / sign convention drift; mitigated by
the fact that Herr 2014 explicitly re-cites Eqs.~1--7 and its numerical
static-limit reproduces our numbers.
\item \textbf{No PIDE-scale $\alpha,\beta$ compilation.} The
$\beta$-vs-$\alpha$ INCONCLUSIVE reflects a coverage limit, not a refutation.
The paper's claim needs the full $\sim$150-line PIDE dataset;
we had 17 lines.
\item \textbf{No radiation-quality (ion / LET) exercise.} RR2964 is
photon-only, but the paper is the entry point to the ion-extended GLOBLE
line. We stayed within its photon-only scope.
\end{enumerate}

\section{Honest assessment of PARTIAL verdict}

This is a textbook LUCID PARTIAL. The self-verdict, the 3-judge audit, and
this backfill all agree: the analytic replication is clean and the
mechanistic dose-response and LQ-map claims are genuinely reproduced, but
the paper's most cited empirical claim (the cross-cell-line
$\beta$-vs-$\alpha$ anti-correlation, Fig.~5) was tested only on a
coverage-limited subset (17 vs $\sim$150 cell lines) and came out with the
wrong sign on that subset. That is a real coverage gap, not a refutation
of the paper; but calling it REPLICATED would over-claim.

Two independent structural gaps also justify PARTIAL rather than
REPLICATED: (i) no track-structure MC was rerun to derive
$\alpha_{DSB}$, so the ``upstream'' half of the GLDM story is
input-not-output here; (ii) no explicit chromatin geometry / DSB-clustering
algorithm was reimplemented, so the ``Poisson-into-loops'' assumption is
inherited whole from the paper rather than stress-tested.

\section{Follow-ups}

See \verb|report/open_questions.json| and
\verb|report/open_questions_section.tex| for the five concrete open
questions with next-step probes (Hi-C-anchored loop distributions,
DSB-repair pathway integration, FLASH-radiotherapy application,
loop-length-distribution sensitivity analysis, ML-driven cluster-severity
scoring). All are free-endpoint feasible.

\section*{Data and code availability}

Code: \verb|code/globle_static.py|, \verb|code/make_figures.py| in this
replication slot (MIT-style clean-room, no author code exists).
Figures: \verb|figures/fig1_dose_response_RT112.png|,
\verb|figures/fig2_alpha_beta_anticorr.png|,
\verb|figures/fig3_decomposition.png|.
No HPC, no GPU, no paid endpoints. Full run reproducible in $<5$~s on a
laptop with \verb|numpy|, \verb|matplotlib|, \verb|scipy|.

\end{document}
